\chapter{Dimensionality Reduction and Multi-scale Representations}

\section{Motivation}

The density fields generated by Direct Numerical Simulations of the Rayleigh--Taylor instability contain a large amount of spatial information distributed across multiple scales. As discussed in the previous chapter, the instability evolves from relatively smooth interface perturbations toward highly complex turbulent structures characterized by sharp gradients, coherent vortices, and intermittent mixing regions.

From a machine learning perspective, directly learning probability distributions in such high-dimensional spaces is particularly challenging. The dimensionality of the data increases rapidly with the spatial resolution of the simulations, leading to high memory requirements and expensive training procedures. Moreover, the limited number of DNS samples available further complicates the learning process and increases the risk of overfitting.

Reducing the dimensional complexity of the data while preserving the physically relevant structures therefore constitutes an important objective. Beyond computational considerations, appropriate representations may also improve the stability and efficiency of generative models by emphasizing dominant structures and reducing redundant information.

Another important aspect concerns the strongly multi-scale nature of RTI flows. Large coherent structures coexist with fine turbulent features and localized interface deformations. Consequently, representations capable of simultaneously capturing both large-scale organization and localized high-frequency structures are particularly attractive.

Several approaches can be considered for this purpose. Classical spectral decompositions based on Fourier transforms provide efficient global frequency representations and are widely used in turbulence analysis. However, they remain limited when dealing with localized or strongly discontinuous structures. Multi-resolution methods based on wavelet transforms constitute an alternative approach capable of jointly representing spatial localization and scale-dependent information.

The objective of this chapter is therefore to introduce the different representations investigated in this work and to analyze their respective advantages and limitations in the context of Rayleigh--Taylor instability fields.

\section{Fourier Representation}

\subsection{Spectral Decomposition}

Fourier transforms constitute one of the most widely used tools for the analysis of turbulent flows and physical fields. The main idea consists in representing a signal as a superposition of sinusoidal modes associated with different spatial frequencies.

For a two-dimensional scalar field $f(x,y)$, the Fourier transform is defined as

\[
\hat{f}(k_x,k_y)
=
\int\int
f(x,y)
e^{-i(k_x x + k_y y)}
\,dx\,dy,
\]

where $(k_x,k_y)$ denote the spatial frequency components.

\begin{figure}[htbp]
\centering
\fbox{
\parbox[c][6cm][c]{0.8\textwidth}{
\centering
Placeholder for Fourier decomposition illustration
}}
\caption{Illustration of a Fourier decomposition of a two-dimensional density field.}
\label{fig:fourier_decomposition}
\end{figure}

In turbulence studies, Fourier analysis is particularly useful because it naturally describes the distribution of energy across spatial scales. Low-frequency modes correspond to large coherent structures, while high-frequency modes represent fine turbulent fluctuations.

\subsection{Advantages for Turbulence Analysis}

Spectral methods possess several important advantages for the analysis of fluid flows. First, Fourier transforms provide compact representations for smooth periodic signals. Second, many physical quantities of interest, such as energy spectra, can be directly expressed in spectral space.

In the context of DNS, pseudo-spectral solvers themselves are often formulated using Fourier representations because of their high numerical accuracy and low dissipation properties. Consequently, Fourier analysis naturally appears as a first candidate for dimensionality reduction and feature extraction.

Furthermore, spectral filtering techniques make it possible to separate structures according to their characteristic scales, which is particularly relevant for turbulent flows exhibiting cascade mechanisms.

\subsection{Limitations for Localized Structures}

Despite their numerous advantages, Fourier representations also exhibit important limitations when dealing with localized or non-smooth structures. Since Fourier basis functions extend over the entire spatial domain, local modifications of the signal affect all spectral coefficients simultaneously. As a consequence, sharp gradients and localized discontinuities generally require a large number of Fourier modes to be accurately reconstructed.

This limitation is closely related to the Gibbs phenomenon, which appears when truncated Fourier series are used to approximate signals containing discontinuities or strong localized variations. Near sharp interfaces, oscillatory artifacts emerge and persist even when increasing the number of retained modes.

\begin{figure}[htbp]
\centering
\fbox{
\parbox[c][6cm][c]{0.8\textwidth}{
\centering
Placeholder for Gibbs phenomenon illustration
}}
\caption{Illustration of the Gibbs phenomenon near a discontinuity reconstructed using a truncated Fourier representation.}
\label{fig:gibbs}
\end{figure}

Although RTI density fields are not strictly discontinuous, they contain steep density gradients, localized interface deformations, and intermittent turbulent structures that exhibit similar reconstruction difficulties. Capturing such localized features using Fourier modes may therefore require a large number of coefficients, reducing compression efficiency and increasing the complexity of the representation.

In addition, Fourier decompositions provide limited spatial localization information. While they accurately describe the global frequency content of the flow, they do not directly indicate where specific structures are located in physical space. This limitation becomes particularly important for turbulent RTI fields, where coherent structures and mixing regions evolve dynamically and remain strongly localized.

These observations motivate the investigation of multi-resolution approaches capable of simultaneously representing both scale-dependent information and spatial localization properties.


\section{Wavelet Transform}

\subsection{Multi-resolution Analysis}

Wavelet transforms provide a representation framework capable of simultaneously capturing spatial localization and scale-dependent information. Unlike Fourier modes, wavelets are localized both in physical space and frequency space.

The central idea of wavelet analysis consists in decomposing a signal onto a family of basis functions generated through scaling and translation operations applied to a mother wavelet.

A continuous wavelet decomposition can be written as

\[
W(a,b)
=
\int
f(x)
\psi_{a,b}(x)
\,dx,
\]

where $a$ represents the scale parameter and $b$ the spatial localization parameter.

\begin{figure}[htbp]
\centering
\fbox{
\parbox[c][6cm][c]{0.8\textwidth}{
\centering
Placeholder for wavelet localization illustration
}}
\caption{Illustration of spatial and frequency localization properties of wavelets.}
\label{fig:wavelet_localization}
\end{figure}

Wavelet decompositions therefore naturally provide hierarchical multi-scale representations in which coarse structures and fine details are separated across different resolution levels.

\subsection{Discrete Wavelet Transform}

In practical applications, discrete wavelet transforms (DWT) are generally preferred because of their computational efficiency. The discrete transform recursively decomposes the signal into approximation and detail coefficients corresponding to different frequency bands.

For two-dimensional fields, the decomposition separates horizontal, vertical, and diagonal structures at multiple scales.

\begin{figure}[htbp]
\centering
\fbox{
\parbox[c][6cm][c]{0.85\textwidth}{
\centering
Placeholder for 2D wavelet decomposition scheme
}}
\caption{Example of multi-resolution decomposition obtained using a two-dimensional discrete wavelet transform.}
\label{fig:wavelet_2d}
\end{figure}

This hierarchical structure is particularly attractive for RTI datasets because turbulent flows naturally contain structures distributed across a wide range of scales.

\subsection{Compression and Sparsity}

One important property of wavelet representations is their ability to produce sparse decompositions for signals containing localized structures or coherent patterns. A large portion of the signal energy can often be concentrated into a relatively small number of significant coefficients.

This property has been extensively exploited in image compression applications such as JPEG2000, where wavelet decompositions provide efficient representations of natural images while preserving localized features.

In the context of RTI simulations, wavelet sparsity may help reduce the effective dimensionality of the dataset while preserving physically important structures such as interfaces, vortices, and localized mixing regions.

Furthermore, the hierarchical organization of wavelet coefficients appears particularly compatible with the multi-scale nature of turbulence and may therefore provide advantageous representations for generative modeling.

\section{Comparison of Reduction Methods}

\subsection{Comparison Criteria}

The different representations introduced in this chapter possess complementary properties. In order to evaluate their relevance for RTI density fields, several comparison criteria can be considered:
\begin{itemize}
    \item reconstruction accuracy,
    \item compression efficiency,
    \item preservation of coherent structures,
    \item localization capabilities,
    \item computational cost,
    \item compatibility with machine learning models.
\end{itemize}

\subsection{Fourier versus Wavelets}

Fourier representations are particularly efficient for globally smooth and periodic structures. They provide natural descriptions of spectral energy distributions and remain highly relevant for turbulence analysis.

However, because Fourier modes are globally supported, localized structures generally require many coefficients to be accurately reconstructed. This limitation becomes important for RTI density fields, which exhibit intermittent structures and sharp localized gradients.

Wavelet transforms, on the other hand, provide localized multi-resolution representations capable of simultaneously capturing large-scale organization and fine-scale structures. Their sparse representations may therefore constitute a significant advantage in the context of limited datasets and generative modeling.

\begin{figure}[htbp]
\centering
\fbox{
\parbox[c][6cm][c]{0.85\textwidth}{
\centering
Placeholder for reconstruction comparison figure
}}
\caption{Comparison of reconstruction quality between Fourier and wavelet-based representations.}
\label{fig:comparison_reconstruction}
\end{figure}

These observations motivate the investigation of wavelet-based diffusion models presented in the following chapters.